% Riemann Framework
% Copyright (C) 2026 MILAN NIKOLIC
%
% This program is free software: you can redistribute it and/or modify
% it under the terms of the GNU General Public License as published by
% the Free Software Foundation, either version 3 of the License, or
% (at your option) any later version.
%
% This program is distributed in the hope that it will be useful,
% but WITHOUT ANY WARRANTY; without even the implied warranty of
% MERCHANTABILITY or FITNESS FOR A PARTICULAR PURPOSE.  See the
% GNU General Public License for more details.
%
% You should have received a copy of the GNU General Public License
% along with this program.  If not, see <https://www.gnu.org/licenses/>.
%
% Build: pdflatex -interaction=nonstopmode -halt-on-error evaluation_protocol.tex
% (run twice for the table of contents and cross-references)

\documentclass[11pt,a4paper]{article}

\usepackage{amsmath,amssymb,amsthm}
\usepackage{booktabs}
\usepackage{geometry}
\usepackage{hyperref}
\usepackage{tabularx}
\usepackage{array}

\geometry{margin=1in}
\setlength{\emergencystretch}{2em}
\setlength{\tabcolsep}{4pt}
\hypersetup{
  colorlinks=true,
  linkcolor=black,
  citecolor=black,
  urlcolor=black,
  pdftitle={A Screening Protocol for Claimed Proofs of the Riemann Hypothesis},
  pdfauthor={Milan Nikolic}
}

\newtheorem{proposition}{Proposition}
\newtheorem{definition}{Definition}
\newtheorem{remark}{Remark}

\newcommand{\RH}{\mathrm{RH}}
\newcommand{\GRH}{\mathrm{GRH}}
\newcommand{\DHf}{\mathrm{DH}}
\newcommand{\Z}{\mathbb{Z}}
\newcommand{\Q}{\mathbb{Q}}
\newcommand{\R}{\mathbb{R}}
\newcommand{\C}{\mathbb{C}}
\newcommand{\verdict}[1]{\textsf{#1}}

\title{\bfseries A Screening Protocol for Claimed Proofs of the\\ Riemann Hypothesis:\\[2pt]
Instrument Domains, Evidence Classes, and an\\ Evaluation of the 2026 Corpus}

\author{Milan Nikoli\'c\\[2pt]
\small Independent Researcher\\
\small \texttt{m\_nikolic@yahoo.com}}

\date{22 November 2026}

\begin{document}

\maketitle

\begin{abstract}
\noindent
Claims of a proof of the Riemann hypothesis are now produced faster than they can be read,
and the instruments available to check them answer different questions: a proof assistant
checks a term, a numerical experiment checks an instance, a logical screen checks a class.
Conflating the three produces verdicts that look decisive and are not. This paper makes the
distinction operational. We give a protocol that routes a claimed proof to the instruments
whose input domain it actually occupies, attaches an evidence class to every verdict, and
uses a six-verdict vocabulary that contains no word meaning ``true''. We then apply the protocol
to a defined, dated, hash-archived corpus: the 2026 Alp\"oge Lean release and its objection,
the 2026 Masullo typescript, the two proportion theorems of Alp\"oge--Furman and Lamzouri, and
the eleven work packages of the screening framework itself. Every verdict is reported with the
evidence it rests on and with the single observation that would overturn it. The framework
rejects one of its own proposals outright and returns ``outside the domain'' for the two
largest claims in the corpus, which is the result the protocol is designed to make visible
rather than to hide. The Riemann hypothesis remains open.
\end{abstract}

\tableofcontents

\section{Introduction}

\subsection{The verification gap}

Three kinds of check are routinely described as ``verifying a proof'', and they establish
different things.

A \emph{kernel} check establishes that a term inhabits a type. It is decisive about the term
and silent about the definitions the term is written against. When the type is Mathlib's own
\texttt{RiemannHypothesis}, a kernel check establishes that the statement is the intended one
and that the proof term has it; that is an achievement, and it is not the same achievement as establishing
that the argument's internal objects carry the mathematics their names suggest.

A \emph{reading} establishes that a human can follow an argument from stated hypotheses to a
stated conclusion. It is decisive about the argument as written and silent about whether the
argument, transcribed, still says the same thing.

A \emph{logical screen} establishes that an argument does not belong to a class already known
to fail. It is decisive in one direction only: it can reject, and it cannot confirm.

The 2026 claims make the gap concrete. Alp\"oge released a Lean archive and no prose proof,
saying that it would be ``obviously inappropriate to attempt to post some massive unreadable
travesty of `slop'\,'', so that ``I decided to only present Lean'' \cite{alpoge}. A Lean and
Mathlib joint statement certified the statement and the term \cite{leanmathlib}; a verification
page recorded the audit trail \cite{sutherland}. Within a day, Nielsen argued that the archive
discharges its geometric obligations vacuously \cite{nielsen}. Both positions can be true at
once, because the kernel check and the objection are about different objects: the first about
a term, the second about what the term's definitions mean. No amount of re-running the first
answers the second, and no amount of reading the second contradicts the first.

\subsection{What an evaluation can be}

This paper does not ask whether a claimed proof is \emph{true}. It asks a question that has
an answer: \emph{which instrument has this claim as an input, and what is the strongest thing
that instrument can say?} The answers are drawn from a fixed vocabulary
(Definition~\ref{def:verdicts}): a claim can be \emph{rejected} by a logical screen, declared
\emph{outside the domain} of every instrument held, left \emph{undetermined} with a named
computation that would decide it, or \emph{validated} only in the finite, elementary or
term-level sense that an instrument actually supports. The qualification ``term-level'' is a
warning as much as a result, since a term can inhabit a type whose internal definitions carry no
mathematics.
There is deliberately no verdict ``true'', and no verdict ``promising''.

\subsection{Contributions}

\begin{enumerate}
\item A table of instrument domains (Table~\ref{tab:instruments}): for each instrument in a
  working screening framework, the input it accepts, the conclusion it licenses, the conclusion
  it does not license, and its characteristic failure mode.
\item A protocol (Section~\ref{sec:protocol}) that assigns each claim an evidence class and a
  verdict, with a pre-registered list of observations that would falsify the study itself.
\item An evaluation of a defined corpus (Section~\ref{sec:evaluations}): the 2026 Alp\"oge
  archive and its objection, the 2026 Masullo typescript, the proportion theorems of
  Alp\"oge--Furman and Lamzouri, and the framework's own eleven work packages, the last as a
  control, because a protocol that cannot reject its author's proposal is not a protocol.
\item A calibration of the Davenport--Heilbronn screen (Section~\ref{sec:calibration}): the
  screen returns ``eligible'' for three of the nine candidate records in the framework's own
  archive, which shows that eligibility is cheap and that the screen's value lies entirely in
  what it rejects.
\item A reproducibility appendix: every artefact cited is stored with a SHA-256 hash, every
  number is produced by a committed test, and the commands are one line each.
\end{enumerate}

\subsection{Scope, stated before the results}

``All works of all mathematicians'' is not a corpus. There is no sampling frame for it: the
population is unbounded, unbounded in time, and, decisively, not observable, because most
arguments exist only as prose that no instrument in this paper can take as input. A study that
claimed to evaluate it would be reporting on the works that happened to be machine-readable
while describing itself as reporting on the works. Section~\ref{sec:corpus} therefore fixes an
inclusion rule in advance, and the census it yields is stated with its exclusions.

\section{The instrument panel}
\label{sec:instruments}

\subsection{Domains}

\begin{table}[t]
\centering
\caption{Instruments, their input domains, and what each one does \emph{not} license.
``Yields'' is the strongest conclusion the instrument supports, not a claim about truth.}
\label{tab:instruments}
\footnotesize
\begin{tabularx}{\textwidth}{@{}>{\raggedright\arraybackslash}p{2.55cm}>{\raggedright\arraybackslash}p{3.05cm}>{\raggedright\arraybackslash}X>{\raggedright\arraybackslash}p{3.5cm}@{}}
\toprule
Instrument & Accepts & Yields & Does not yield \\
\midrule
Idea pipeline A--E & a closed-form map or an operator, plus a falsification criterion &
  ``not falsified by any gate it was subjected to'', stage by stage &
  correctness; it returns noise for a proof-shaped claim, which has no map \\
\addlinespace
Davenport--Heilbronn (DH) screen & a set of \emph{declared} inputs, each classified against a
  sourced inventory & \verdict{dead}: every input holds for DH, where the RH-analogue is false,
  so the argument would prove a false statement & anything positive; \verdict{eligible} means
  only that the contrapositive does not fire \\
\addlinespace
Spectral statistics & an eigenvalue spectrum & location and spacing statistics, versus GOE/GUE
  references & anything about the zeros; the zeros are not an input \\
\addlinespace
Lean + axiom audit & a formalisation and a statement & the term inhabits the statement; the
  axiom set; that a recorded open target still contains \texttt{sorry} &
  that the statement means what its name says, or that internal definitions carry content \\
\addlinespace
Finite certificate & an exactly decidable finite statement & a proof of that statement &
  any universal statement; a verified instance is not a verified family \\
\addlinespace
Elementary certificate & an explicit inequality, a constant chain, a tight case & a proof of
  the certificate, which is the paper's load-bearing elementary step & the analytic inputs on
  which the paper's theorem also depends \\
\addlinespace
Archive with hashes & a document & that the bytes quoted are the bytes stored, and that a
  silent revision is detectable & that the document is correct, or even readable \\
\bottomrule
\end{tabularx}
\end{table}

Table~\ref{tab:instruments} is the paper's central object; the protocol in
Section~\ref{sec:protocol} is a procedure for reading it. Two rows deserve commentary because
they are the ones most often mistaken for each other.

\subsection{The DH contrapositive, and why eligibility is cheap}
\label{sec:calibration}

Let $f$ be a Dirichlet series with a functional equation but no Euler product, whose zeros are
known not all to lie on the critical line; the Davenport--Heilbronn function is the standard
example \cite{leclair,vorsos}. Then:

\begin{proposition}[contrapositive screen]
\label{prop:dh}
If every ingredient of an argument holds for $\DHf$, the argument cannot establish $\RH$.
\end{proposition}

\begin{proof}
Apply the argument to $\DHf$. Its conclusion, transferred verbatim, asserts the analogue of
$\RH$ for $\DHf$, which is false. A valid argument with true premisses cannot have a false
conclusion.
\end{proof}

The screen is only as good as its inventory, and the inventory must be sourced rather than
assumed. The framework's implementation classifies nine inputs, each with a citation and, where
possible, a computation: three transfer to $\DHf$ (the functional equation, real coefficients,
zero symmetry), three fail for it (the Euler product, the prime side of the explicit formula,
Li positivity), and three are left unsourced rather than filled in from expectation (the
Hadamard order, the zero-counting asymptotic, GUE pair correlation). Three of the nine are
verified by computation inside the repository, including the single arithmetic fact that
settles the Euler-product row: for the DH coefficients $b$, $b(2)b(3)=-\kappa^{2}\neq b(6)=1$,
so the coefficients are not multiplicative.

Now the calibration. The framework's own archive holds nine candidate records. Screen all nine
and the dead list has \emph{one} entry: the framework's own flagship proposal, which declares
only the functional equation and the zero symmetry, both of which $\DHf$ has. Three records are
\verdict{eligible}, five are \verdict{unscreened}. Eligibility is therefore not scarce: any
argument that touches the prime side of an explicit formula qualifies, which is most of
analytic number theory. The screen's epistemic value is concentrated in the dead list, and a
paper that reported only ``three of nine pass'' would be reporting the wrong statistic.

A second limit is sharper and is stated here because it applies directly to
Section~\ref{sec:masullo}. The screen tests \emph{transfer}, not \emph{circularity}. An
argument that assumes an $\RH$-equivalent hypothesis (Li's criterion, say) is screened as
\verdict{eligible} for DH, because Li's criterion fails for DH. The screen cannot tell
discrimination from circularity. Passing it is necessary and nowhere near sufficient.

\subsection{Finite and universal}

The framework contains two proofs that a computation \emph{is} a proof, because the statements
are finite and exactly decidable: a certified sieve establishes $p_{1{,}000{,}000}=15{,}485{,}863$
with no gap below it, and exact integer arithmetic establishes the Riemann hypothesis for
individual curves over finite fields: 942 elliptic curves over $\mathbb{F}_{5},\dots,
\mathbb{F}_{19}$, 324 smooth genus-$2$ curves over $\mathbb{F}_{3}$, and 10{,}000 over
$\mathbb{F}_{5}$, with zero failures. Neither is a numerical sample; both are finite claims
decided exactly.

The hypothesis over $\Q$ is universal, and no finite computation establishes a universal
statement. This is not a limitation of effort. It is the reason the rows of
Table~\ref{tab:instruments} are not interchangeable, and the reason a ``verified to height
$T$'' result is evidence about the height and not about the zeros.

\subsection{What formal verification can and cannot see}

The 2026 archive raises the point in its sharpest form. A kernel replay, an axiom audit and a
statement comparison answer, respectively: does the term typecheck; which axioms does it use;
is the theorem the intended statement. None of the three answers the question whether an
internal definition says what its name says. If a predicate is defined as
\texttt{fun \_ \_ \_ => True}, then a certificate built from it is a valid term, and every
downstream step that consumes the certificate typechecks. Formally:

\begin{proposition}[vacuous discharge]
\label{prop:vacuous}
Let $C$ be an inductive closure with a constructor accepting any proof of a predicate $\Gamma$.
If $\Gamma$ is constantly true, then $C\,\Gamma$ holds at every index, and the proof has used no
hypothesis.
\end{proposition}

\begin{proof}
Immediate by the constructor and $\beta$-reduction: the constructor's argument is a proof of
$\Gamma\,x$, which is definitionally $\top$; \texttt{trivial} supplies it.
\end{proof}

Proposition~\ref{prop:vacuous} is machine-checked in the framework
(\texttt{Dispute.lean}, sorry-free for this fragment), and it is the
checkable core of Nielsen's objection. It is also, by itself, not an objection to any
particular archive: an archive may contain vacuous predicates that no proof calls. The decisive
question is a dependency question, and it is not a kernel question.

\subsection{The evidence chain}

Every document this paper evaluates is stored byte for byte with a SHA-256 hash and a page
count, in the repository's \texttt{references/} directory. This is not bookkeeping. Two of the
three documents are being actively revised by their authors, and a hash is what makes a silent
revision detectable; the third is served from a personal site and from an archive that answers
automated retrieval with a challenge, so a citation without stored bytes is a citation no
future reader can check.

\section{Protocol}
\label{sec:protocol}

\subsection{Evidence classes}

\begin{definition}[evidence class]
For a claim $c$, write $E(c)\subseteq\{0,1,2,3,4\}$ for the classes it attains:
$E0$: the primary document exists in a form that can be archived and hashed;
$E1$: a statement can be extracted from it, with a location;
$E2$: its declared inputs can be identified and classified;
$E3$: a machine-checkable artefact accompanies it (a formalisation, a certificate, a
reproducible computation);
$E4$: an independent party has reproduced a check on it.
\end{definition}

The classes are nested in cost, not in strength. A claim at $E3$ with no $E1$ (a formalisation
whose theorem is not stated in prose) is not better off than one at $E1$; it is differently
situated, and the protocol's job is to say how.

\subsection{Decision procedure}

Given a claim $c$:
\begin{enumerate}
\item \textbf{Archive} the primary document; record hash and length. If it cannot be archived,
  stop: the claim is not evaluable here, and say so.
\item \textbf{Extract the statement}, with page or line. If there is no statement (a
  placeholder, an unfinished chapter), record that fact; it is a result about the document and
  it bounds everything that follows.
\item \textbf{Identify the declared inputs} and classify them. If they are enough to run the DH
  screen and the screen fires, the claim is \verdict{rejected}; stop.
\item \textbf{Route to an instrument} whose input domain the claim occupies: a map for the
  pipeline, a spectrum for the statistics, a formalisation for the kernel, a finite object for a
  certificate. If none accepts it, the claim is \verdict{outside domain}.
\item \textbf{Record the verdict}, its evidence class, and the single observation that would
  change it.
\end{enumerate}

\subsection{Verdicts}
\label{def:verdicts}

\begin{definition}[verdict vocabulary]
\verdict{rejected}: a logical screen applies and fires. \verdict{outside domain}: no instrument
held accepts the claim as input. \verdict{undetermined}: an instrument applies and has not been
run to completion, or its result is contested; the decisive computation is named.
\verdict{validated (finite)}: an exactly decidable finite statement has been proved.
\verdict{validated (elementary)}: the load-bearing elementary components of a theorem have been
proved, with the analytic inputs named and outside the validation.
\verdict{validated (term)}: a proof term inhabits a stated type and uses only the standard
axioms; this says nothing about whether the archive's internal definitions carry the
mathematics their names suggest. No verdict asserts that a claim is true, and the absence of
\verdict{rejected} is not evidence for a claim.
\end{definition}

\subsection{Pre-registered falsification criteria for this study}

In the sense of Popper, a study that cannot fail is not a study. This paper is refuted by any
one of the following:
\begin{enumerate}
\item a claim recorded as \verdict{rejected} whose declared-input set is shown to include an
  input that fails for DH, i.e.\ a misclassification in the inventory;
\item a claim recorded as \verdict{outside domain} that is shown to occupy a domain of an
  instrument held;
\item a document whose archived hash does not match the hash quoted here;
\item a number in Section~\ref{sec:evaluations} that the committed tests do not reproduce.
\end{enumerate}

\section{The corpus}
\label{sec:corpus}

\subsection{Why ``all works'' is not a corpus}

A census needs a frame: a rule that decides membership before the results are seen, and that
could in principle be applied to every candidate. For claims about $\RH$ there is no such rule
at the scale of the literature, for three reasons. The population is open: new claims appear
daily. The claim type is heterogeneous: a claim may be a paper, a formalisation, a lecture, or
a file; the instruments here accept only some of these. And the observable subset is
unrepresentative by construction: an argument with no formalisation and no certificate cannot
be screened by anything in this paper, so a census restricted to screenable claims measures the
availability of artefacts, not the state of the subject.

The honest object is therefore a \emph{defined corpus with a stated inclusion rule}, and the
results are about that corpus. The rule used here:

\begin{quote}
A work is in the corpus if (i) a primary document is available for archival, (ii) it was
released on or after 20 November 2026 and is publicly attributed to a claimed proof of $\RH$ or
$\GRH$ or to a check of such a claim, and (iii) its declared inputs or its finite content can be
identified from the document itself.
\end{quote}

Clause (ii) is a date. It is there because the framework's evaluation machinery is itself dated,
and comparing it to arguments written before the machinery existed would confound the two.

\subsection{The census}

\begin{table}[t]
\centering
\caption{The corpus. Hashes are SHA-256 of the stored file; page counts are from the PDF page
tree. The framework rows are work packages rather than external claims.}
\label{tab:corpus}
\footnotesize
\begin{tabularx}{\textwidth}{@{}>{\raggedright\arraybackslash}p{2.7cm}>{\raggedright\arraybackslash}p{2.2cm}>{\raggedright\arraybackslash}p{2.45cm}>{\raggedright\arraybackslash}X@{}}
\toprule
Work & Date & Extent & Archived / recorded as \\
\midrule
Alp\"oge, \emph{Remark} & 2026-11-20 13:52 & 4 pp. &
  \texttt{alpoge\_remark.pdf}, hash \texttt{3fc97784\ldots} \\
Lean archive + joint statement + verification page & 2026-11-20/21 & 175{,}187 modules &
  cited, not stored (61.5~GB); statement and audit at \cite{leanmathlib,sutherland} \\
Nielsen, \emph{Invalidity} & 2026-11-20, PDF 11-21 16:50 & 13 pp. &
  \texttt{nielsen\_invalidity.pdf}, hash \texttt{2f8bd3c4\ldots} \\
Masullo, \emph{The Music of the Primes} & 2026-11-20 19:13 & 940 pp. &
  \texttt{masullo\_music.pdf}, hash \texttt{e039d3f6\ldots} \\
Alp\"oge--Furman & 2026 & arXiv:2608.13637 &
  certificates in \texttt{proportion\_certificates.py} \\
Lamzouri & 2026 & preprint & arXiv:2609.02882; the same module \\
Framework: eleven work packages & 2026 & 11 Lean modules, 1272 tests &
  \texttt{lean/}, \texttt{python/}, \texttt{ideas/} \\
\bottomrule
\end{tabularx}
\end{table}

\subsection{Exclusions}

Excluded, with what their inclusion would require: arguments available only as prose without a
PDF (needs transcription and archival, and then only the DH screen applies); claims about
$\RH$ over function fields (out of scope: $\RH$ is a theorem there, so the analogue of the
question is different); and the 2026 PhilArchive critique of the Masullo draft, which is
reported but was not retrieved; it is named in the repository's record and is not counted as
evidence anywhere in this paper.

\section{Evaluations}
\label{sec:evaluations}

\subsection{The Alp\"oge archive: a kernel verdict on a term}

\textbf{Evidence.} The primary prose is four pages \cite{alpoge}: a statement of belief, a
20-entry timeline, and an explanation of the medium. The archive itself is 17.4 million lines
in 175{,}187 modules, with a stated dependency on Mathlib alone. The Lean and Mathlib
communities certified the statement and the term \cite{leanmathlib}; Sutherland's page records
the mechanical audit \cite{sutherland}. As of the date of the record, the audit reports
175{,}187 modules compiling, axioms exactly \texttt{propext}, \texttt{Classical.choice},
\texttt{Quot.sound}, and Mathlib's compiled artefacts byte-identical to the pin, with the
comparator run and the paranoid-kernel replay still outstanding at the time the record was
written.

\textbf{What the instruments say.} The statement is Mathlib's \texttt{RiemannHypothesis}; the
term has that type; the axiom set is the standard one. In the vocabulary of
Definition~\ref{def:verdicts}, that is a \verdict{validated} verdict about a \emph{term}, and
that verdict is not available for reuse as a verdict about the \emph{argument}, because no
argument is written down. The author says why: a prose write-up would be a ``massive unreadable
travesty of `slop'\,'' \cite{alpoge}.

\textbf{The objection, and its status.} Nielsen's 13-page reading of the source \cite{nielsen}
argues that the hypothesis bundle is inhabited through vacuous witnesses: a geometry predicate
defined as \texttt{fun \_ \_ \_ => True}; a degenerate model whose positivity fields are
$\top$ and whose index types are empty; a map from lattices to sections that is the zero map; a
lattice identification that holds by \texttt{rfl}; and a height equation that holds because the
height was assigned to be its own formula. Her Section~8 collects the findings; her Section~9
lists four checks that would settle the matter, of which the decisive one is a transitive
dependency trace from the final theorem to the definitions that discharge the geometric
obligations; and her Section~10 states the limits of her own method: about 70 files read of
175{,}187, no build, no axiom print, no compiled trace.

The protocol's verdict on the objection is \verdict{undetermined}, with the decisive
computation named: the dependency trace. The protocol's verdict on the certification is
\verdict{validated (term)} together with the observation that a term-verdict is silent about
internal definitions. This is the paper's central worked example of the gap in
Section~1.1, and it is why the framework stores the logical core of the objection as a proved
lemma (Proposition~\ref{prop:vacuous}) while refusing to state the objection itself as a
theorem about \texttt{RiemannHypothesis}: a claim about source code cannot be expressed as a
Lean proposition about $\RH$, and pretending otherwise would be an equivocation.

\textbf{One open, decidable item.} Nielsen names four identifiers for the degenerate model
(\texttt{geoDegenerate}, \texttt{algDegenerate}, \texttt{padDegenerate},
\texttt{mdlJunk}) as ``the release's own''. A third-party audit reports that every
\texttt{degenerate} hit in the tree is the ordinary word \texttt{nondegenerate} \cite{chu}.
One of the two is wrong, and the test is a single case-sensitive search over the archive. The
framework records the conflict and does not resolve it: it does not hold the 61.5~GB export. A
case-sensitive search for \texttt{degenerate} would not match \texttt{geoDegenerate}, which is
a plausible mechanism for the discrepancy; that hypothesis is itself checkable and is not
checked here.

\subsection{The Masullo draft: no statement to screen}
\label{sec:masullo}

\textbf{Evidence.} A 940-page typescript \cite{masullo}, archived with its hash. Its
acknowledgements describe a draft: ``the math is complete, and only filling in some writing
remains'', ``the introduction itself is unfinished'', and ``several glitches here and there that
I was in the process of addressing''. The file was created on 2026-11-20 at 19:13, six hours
after Alp\"oge's remark.

\textbf{The statement is a placeholder.} Chapter~13, ``Zeros of Hecke $L$-functions'', occupies
pp.~521--522. Its Section~13.3 is titled \emph{Main Theorems on the location of the zeros}.
Theorem~13.3.1.1 (Trace Formula) is stated. The text then reads ``Upon combining Theorem
13.3.1.1 and Theorem ??, we obtain the following:''. Theorem~13.3.1.2, \emph{Divisor of
$L(s,\chi)$}, is printed with its number and title and no statement at all; the footnotes on
that page read ``Correct statement.'', ``Insert statement.'', and ``State and prove the Main
Theorems, also explaining how to deduce them by approximation
$\varphi\to e^{-s|\log t|}$ as $\mathrm{Re}(s)>1$.'' The GRH content is attributed to that
theorem one chapter later, at p.~528, under the author's own footnote ``Correct this
subsection.''

In the protocol's terms this is a finding at step~2, and it bounds the evaluation: a theorem
that is not stated is not screened, and a verdict about it would be a verdict about a name.
The correct output is \verdict{outside domain} for every instrument except one.

\textbf{The one instrument that applies.} The draft's declared inputs can be read off its own
sections: the Euler product (\S1.4.7, \S14.1.2, where partial Euler products are used); the
prime side of Weil's explicit formula (\S3.4, Appendix B.4; the stated trace formula carries an
explicit sum over finite places weighted by $\log q_v$); and the functional equation
(\S14.1.1, Appendix B.4). Run the DH screen on that declaration and the verdict is
\verdict{eligible}: the two arithmetic inputs fail for DH, so the contrapositive of
Proposition~\ref{prop:dh} does not fire. The screen's calibration
(Section~\ref{sec:calibration}) applies with full force: three of nine records in the
framework's archive are also eligible, and eligibility is permission, not promise.

Two absences from the declaration are as informative as its contents. The draft says at
\S1.4.6 that Weil's positivity criterion is \emph{not} used; declared, it would have misstated
the paper and would not have strengthened the verdict, since Li's criterion fails for DH and the
screen would have read \verdict{eligible} from an $\RH$-equivalent input. And the trace formula
is proved only for test functions with $\varphi(C_F^{1})=0$, a restriction the text says twice
``cannot be removed before GRH is proved'' (pp.~30, 35). The removal step is exactly what the
missing Chapter~13 would supply. The draft therefore describes a strategy whose two load-bearing
steps are, on its own account, not yet carried out.

\textbf{Verdicts.} \verdict{outside domain} for the numerical instruments (there is no
spectrum), for the kernel (there is no formalisation), and for the mathematics as posted,
because the theorem is not stated. Once Theorem~13.3.1.2 is stated and proved the verdict
becomes \verdict{undetermined}, with the decisive observations already named: the statement
itself, and the removal of the test-function restriction. The DH screen returns
\verdict{eligible} for the declared inputs, and the protocol records that this is the weakest
positive label in its vocabulary.

\subsection{The proportion theorems: what an elementary certificate can and cannot do}

\textbf{Evidence.} Alp\"oge--Furman \cite{alpogefurman} prove that more than two thirds of the
non-trivial zeros are simple and on the critical line ($0.6725$ with the Montgomery--Taylor
window), and Lamzouri \cite{lamzouri} proves the same for more than five sixths ($0.83625$
distinct). Both are proportion theorems and both say so.

\textbf{The certificates.} The framework implements the elementary parts: a rank--trace
inequality, a multiset inequality for conjugation-invariant multisets, the constant chain
$2-R(\psi)$ and $(3-R(\psi))/2$, and the admissibility of the kernel used. The tests weight the
explicit tight cases above random search, because a Monte Carlo search cannot prove an
inequality; it can only fail to find a counterexample. Measured: the rank--trace inequality
held in 20{,}000 random admissible cases with worst margin $0.096$; the multiset inequality
held in 3{,}000 conjugation-invariant multisets with tightest margin $0.0$: the tight case
again, which is what pins the normalisation; and the constants evaluate to
$0.6725007037$ and $0.8362503518$, matching both papers.

\textbf{Verdict.} \verdict{validated (elementary)} for those components, with the analytic
inputs named and outside the validation: Weil's explicit formula, the unconditional prime-side
second moment, Montgomery--Vaughan, Chebyshev--Mertens, Stirling. And a \verdict{rejected}
verdict for the stronger reading, on the papers' own evidence: their inputs hold for
Davenport--Heilbronn and Epstein zeta functions, for which the analogue of $\RH$ is false, so
by Proposition~\ref{prop:dh} the method cannot reach $\RH$. That is not a caveat attached to
the result; it is a proof that this method has a ceiling, and it is the cleanest example in the
corpus of a paper being right about what it does not do.

\subsection{The framework's own work packages: the control}

\begin{table}[t]
\centering
\caption{Self-application. The framework's own proposals are screened with the same protocol;
one is rejected by the screen it was built to survive, and the programme is abandoned by its own
pre-registered criteria.}
\label{tab:self}
\footnotesize
\begin{tabularx}{\textwidth}{@{}>{\raggedright\arraybackslash}p{4.2cm}>{\raggedright\arraybackslash}p{2.4cm}>{\raggedright\arraybackslash}X@{}}
\toprule
Work package & Verdict & Evidence \\
\midrule
\texttt{dimension\_shift\_w} (flagship) & \verdict{rejected} &
  declares only the functional equation and zero symmetry; both hold for DH; one dead entry in
  the nine-record archive \\
primonic lift of the same operator & \verdict{undetermined} &
  declares the Euler product, which DH lacks; one of the three eligible records \\
DSH sector programme & \verdict{rejected} (by its own criteria) &
  abandoned: GUE input blocks give 2/2 GUE sectors, GOE blocks 0/2, so the class is inherited
  from the input, not produced by the involution; criterion (1) holds only class-conditionally \\
shift-zeta graded lift & \verdict{rejected} &
  the lifted Euler product does not constrain the zeros of $\zeta$; recorded as a negative
  result rather than repaired \\
DSIN / BB84 layer & \verdict{rejected} &
  closed on an information-theoretic obstruction; retained as an instrument \\
finite-field curves & \verdict{validated (finite)} & 942 elliptic over
  $\mathbb{F}_5$--$\mathbb{F}_{19}$, 324 genus-2 over $\mathbb{F}_3$, 10{,}000 over
  $\mathbb{F}_5$, exact integer arithmetic, zero failures \\
prime certificate & \verdict{validated (finite)} & $p_{1{,}000{,}000}=15{,}485{,}863$ with a
  sieve invariant forbidding a gap below it \\
elementary certificates & \verdict{validated (elementary)} & 20{,}000 and 3{,}000 random cases
  plus the tight cases; constants to ten digits \\
\bottomrule
\end{tabularx}
\end{table}

Section~\ref{sec:evaluations} would be worthless if the protocol could only reject other
people's work. Table~\ref{tab:self} is the control. The framework's flagship proposal is the
single entry on the dead list; its sector programme was abandoned by criteria fixed before the
measurements; and its negative results are retained as findings rather than quietly dropped.
The same instrument that returns \verdict{eligible} for the Masullo draft returns
\verdict{rejected} for the framework's own operator, which is the only calibration of a screen
that matters.

\subsection{Results}

\begin{table}[t]
\centering
\caption{All verdicts, with the evidence class and the observation that would overturn the
verdict. ``Validated'' is qualified in every row: finite means a finite statement, elementary
means the elementary components of a theorem.}
\label{tab:results}
\footnotesize
\begin{tabularx}{\textwidth}{@{}>{\raggedright\arraybackslash}p{3.3cm}>{\raggedright\arraybackslash}p{2.0cm}>{\raggedright\arraybackslash}p{1.05cm}>{\raggedright\arraybackslash}X@{}}
\toprule
Work & Verdict & $E$ & Would be overturned by \\
\midrule
Alp\"oge archive (term) & validated (term) & 3,4 & a rejected kernel replay, or a statement
  mismatch \\
Alp\"oge argument (as presented) & outside domain & 0,1 & a prose or machine-readable argument
  that can be screened \\
Nielsen objection & undetermined & 1,2 & the dependency trace, either way \\
Masullo draft (theorems) & outside domain & 0,1 & a stated and proved Theorem 13.3.1.2 \\
Masullo draft (declared inputs) & eligible & 1,2 & a misclassified inventory row \\
Alp\"oge--Furman & validated (elementary) & 3 & a counterexample to Lemma 3.2, or a wrong
  constant \\
Lamzouri & validated (elementary) & 3 & a counterexample to Proposition 2.1 \\
Both proportion methods as $\RH$ routes & rejected & 0,1 & an input that fails for DH \\
Framework flagship operator & rejected & 1,2 & an input that fails for DH \\
DSH programme & rejected & 3 & a sector class not inherited from the input \\
Finite-field / prime certificates & validated (finite) & 3 & an exact counterexample \\
\bottomrule
\end{tabularx}
\end{table}

The distribution is the paper's headline result, and it is deliberately unflattering to the
corpus. Of eleven entries, three are rejected: two of them methods rejected on their own
evidence, one of them the framework's flagship proposal. Two are outside every domain held,
one is undetermined with a named decisive computation, one is eligible without being validated,
and four are validated: one as a statement about a proof term, one as a finite statement, and
two in the elementary sense, with their analytic inputs explicitly outside the validation. Not
one entry is, or could be, a validated proof of $\RH$.

\section{Threats to validity}

\paragraph{Instrument validity.} The DH screen rejects only when every declared input transfers,
and it accepts only declaration sets. An author who omits an input, or names one the inventory
cannot classify, moves the verdict without changing the argument. \emph{Mitigation:} inputs are
declared in a file with the section reference that shows them in use; an unknown key degrades to
\verdict{unscreened} rather than passing; and the record and the module are compared by test so
they cannot drift apart. \emph{Residual:} a deliberate omission is not detectable by the screen.

\paragraph{Screen calibration.} Three of nine records are eligible, so eligibility is weak
evidence. \emph{Mitigation:} the paper reports the dead list, not the eligible list, as the
screen's output. \emph{Residual:} readers may still quote the label.

\paragraph{Corpus bias.} The inclusion rule requires a machine-readable or archival document.
That selects for claims made by people who write formalisations, which is not a random sample of
claims. \emph{Mitigation:} the rule is stated before the results, and the exclusions are listed
with what would be needed to include them. \emph{Residual:} the census cannot be extrapolated.

\paragraph{Evidence-class inflation.} A claim with a formalisation may be treated as better
supported than a claim with a reading, whereas the two are differently supported. \emph{Mitigation:}
the verdict vocabulary has no ordering, and Table~\ref{tab:results} prints the evidence class
next to the verdict rather than folding them together.

\paragraph{Author position.} The framework's author conducts the evaluation, including of his
own proposals. \emph{Mitigation:} the self-application table is a pre-registered control, the
criteria for abandoning the flagship programme were fixed before the measurements, and the
sentences that could be self-serving (``eligible'', ``undetermined'') are the weakest in the
vocabulary, with the strongest reserved for the finite statements. \emph{Residual:} independent
replication is the only real remedy and is not claimed.

\paragraph{Archival dependence.} Hashes protect the bytes, not the availability of the source.
\emph{Mitigation:} the three documents are stored in the repository, not merely linked.

\section{Related work}

The Xena Project's four criteria \cite{xena} are criteria about statements and terms; this paper
adopts them as one row of Table~\ref{tab:instruments} and adds the rows they do not cover. The
Lean and Mathlib statement \cite{leanmathlib} and Sutherland's verification files
\cite{sutherland} supply the mechanical audit whose scope is analysed in \S\ref{sec:instruments}.
Nielsen's analysis \cite{nielsen} is the corpus's clearest instance of a reading-based objection
to a kernel-certified term, and her Section~10 limits are quoted rather than paraphrased away.
The proportion theorems \cite{alpogefurman,lamzouri} are the corpus's clearest instance of a
paper stating its own ceiling; the Davenport--Heilbronn literature \cite{leclair,vorsos,franca}
supplies the counterexample class that makes the ceiling a theorem. The framework's own negative
results are collected in a lessons-learned document that this paper treats as data.

\section{Conclusion}

The instruments answer different questions, and the discipline that matters is not running more
of them but assigning each claim to the one whose input domain it occupies, and saying
``outside domain'' when there is none. Applied to the 2026 corpus, the protocol yields three
rejections, two claims outside every domain, one undetermined entry with a named decisive
computation, one eligible-but-unvalidated declaration, and four validated statements none of
which is a proof of the Riemann hypothesis. The two largest claims in the corpus are the two
that cannot be evaluated: one has no prose argument, the other has no statement yet. That is an
uncomfortable result for a paper about evaluation, and it is the correct one.

The Riemann hypothesis remains open. Nothing in this paper is evidence for or against it, and
no verdict in Table~\ref{tab:results} would change if it were settled tomorrow, except the
\verdict{undetermined} entry, which would stop being interesting.

\appendix
\section{Reproducibility}

All commands run from the repository root; the Python interpreter is the one the project
pins, and every number quoted above is produced by a committed test.

\begin{verbatim}
python -m pytest python/tests -q            # 1272 passed
lake env lean lean/RiemannFramework/Dispute.lean
                                          # two 'sorry' warnings
python scripts/run_dh_filter.py              # dead / eligible lists
python -m pytest python/tests/test_reference_archive.py -q
\end{verbatim}

The paper is itself guarded, on the principle that a document about verification should be
verifiable. \texttt{python/tests/test\_evaluation\_paper.py} extracts this PDF's text layer and
checks four things: that the page and character counts are the compiled article's; that the
sentence keeping the hypothesis open is present; that none of the six phrases asserting a proof
appears; and that the sixteen-hex-digit hashes printed in Table~\ref{tab:hashes} are prefixes of
the SHA-256 digests of the archived files. It cannot check whether the verdicts are right; only
the criteria in \S\ref{sec:protocol} can do that, and they are printed in the paper for a reader
to apply against the author.

\begin{table}[h]
\centering
\caption{Archived documents: hash, length, pages. The index records provenance and the
documents that cite each one.}
\label{tab:hashes}
\footnotesize
\begin{tabular}{@{}lrrl@{}}
\toprule
File & Bytes & Pages & SHA-256 (first 16 hex) \\
\midrule
\texttt{alpoge\_remark\_2026-11-20.pdf} & 245{,}114 & 4 & \texttt{3fc97784b71ddd44} \\
\texttt{nielsen\_invalidity\_2026-11-20.pdf} & 98{,}048 & 13 & \texttt{2f8bd3c47c211457} \\
\texttt{masullo\_music\_of\_the\_primes\_2026-11-20.pdf} & 6{,}756{,}996 & 940 &
  \texttt{e039d3f6cf0a4b14} \\
\bottomrule
\end{tabular}
\end{table}

\paragraph{Declarations.} The author wrote the framework and conducted the evaluation; the
conflict is declared in Section~\ref{sec:evaluations} and mitigated by the control in
Table~\ref{tab:self}. No external funding. The code is GPL-3.0; the archived documents remain
the property of their authors and are stored for verification only.

\begin{thebibliography}{99}

\bibitem{alpoge} L.~Alp\"oge, \emph{Remark}, 20 November 2026. Archived as
  \texttt{references/alpoge\_remark\_2026-11-20.pdf}.

\bibitem{leanmathlib} Lean and Mathlib communities, \emph{Lean verification of the Riemann
  Hypothesis}, 21 November 2026.

\bibitem{sutherland} A.~Sutherland, \emph{Verification files for the Lean proof of the Riemann
  Hypothesis}, MIT, updated 21 November 2026, 12:55 UTC.

\bibitem{nielsen} J.~L.~Nielsen, \emph{The Invalidity of Alp\"oge's Proposed Proof of the
  Riemann Hypothesis}, PhilArchive NIEIOT-2, 20 November 2026.

\bibitem{chu} F.~Chu, \emph{GRH proof's geometry is \texttt{fun \_ \_ \_ => True}, in 24
  files}, 22 November 2026 (third-party audit of the code claims).

\bibitem{masullo} A.~M.~Masullo, \emph{The Music of the Primes}, 940 pp., 20 November 2026.
  Archived as \texttt{references/masullo\_music\_of\_the\_primes\_2026-11-20.pdf}.

\bibitem{alpogefurman} L.~Alp\"oge and R.~Furman, \emph{More than two thirds of the zeros of
  the Riemann zeta function are simple and on the critical line}, arXiv:2608.13637, 2026.

\bibitem{lamzouri} Y.~Lamzouri, \emph{A new proof that more than $2/3$ of the zeros are simple
  and on the critical line}, arXiv:2609.02882, 2026.

\bibitem{leclair} A.~LeClair and G.~Mussardo, on the absence of an Euler product for the
  Davenport--Heilbronn function, as recorded in the framework's DH inventory.

\bibitem{vorsos} A.~Voros, on Li's criterion reacting to the Davenport--Heilbronn zeros, as
  recorded in the framework's DH inventory.

\bibitem{franca} G.~Fran\c{c}a and A.~LeClair, on the functional equation of the
  Davenport--Heilbronn function.

\bibitem{xena} The Xena Project's four criteria for a formalisation (statement, statement
  match, axiom audit, independent replay), as applied in \cite{leanmathlib}.

\end{thebibliography}

\end{document}
